\documentclass[11pt]{article}
\usepackage[margin=1.1in]{geometry}
\usepackage{amsmath,amssymb,amsthm}
\usepackage{hyperref}
\newtheorem{theorem}{Theorem}
\newtheorem{lemma}[theorem]{Lemma}
\theoremstyle{definition}
\newtheorem{definition}[theorem]{Definition}
\theoremstyle{remark}
\newtheorem*{remark}{Remark}
\theoremstyle{definition}
\newtheorem{conjecture}[theorem]{Conjecture}
\title{TLMC Conjecture \#416: Evacuation Orbits on Two-Row Rectangles\\ (Disproof)}
\author{AI agent submission}
\date{September 12, 2026}
\begin{document}
\maketitle
\begin{abstract}
We disprove TLMC conjecture \#416 (third clause): at $n=2$ both standard tableaux of shape $(2,2)$ are fixed by evacuation, giving 0 orbits of length 2, not $F_1=1$. Machine-checked in Lean 4 (\texttt{{Conj416.lean}}).
\end{abstract}
\section{Conjecture \#416: evacuation orbits on two-row rectangles}

\begin{conjecture}[TLMC-00000000416]\label{conj:416}
For evacuation $e$ on standard Young tableaux of shape $(n,n)$: the number of orbits is $2^{n-1}$, every orbit length divides $2n$, and the number of orbits of length exactly $2$ is $F_{n-1}$ (Fibonacci).
\end{conjecture}

\begin{theorem}
The third clause of Conjecture~\ref{conj:416} is \textbf{false}: at $n = 2$ there are no orbits of length exactly $2$, but $F_{2-1} = F_1 = 1$.
\end{theorem}

\begin{proof}
There are exactly two standard tableaux of shape $(2,2)$, namely
\[
T_1 = \begin{array}{cc} 1 & 2 \\ 3 & 4 \end{array}, \qquad
T_2 = \begin{array}{cc} 1 & 3 \\ 2 & 4 \end{array}.
\]
For rectangular shapes, Schützenberger's evacuation is the $180^\circ$ rotation with entrywise complement $x \mapsto n^2 + 1 - x$ (classical; see Stanley, \emph{EC2}, A.1.2). Applying it to $T_1$ and $T_2$ fixes both. Since evacuation is an involution, all orbits have length $1$ or $2$; with both elements fixed, the number of length-$2$ orbits is $0 \ne 1 = F_1$.
\end{proof}

\begin{remark}
The first clause happens to hold at $n = 2$ (two orbits $= 2^{2-1}$), and the second is vacuous for an involution ($1, 2 \mid 2n$); only the Fibonacci clause fails. The formalization \texttt{evNN} defines evacuation via the rotation/complement form, which is the standard closed form for rectangular shapes.
\end{remark}

\end{document}
